\makeatletter
\expandafter\gdef\csname CGTip@sec:introduction\endcsname{Section 1: Introduction (p. 3)}
\expandafter\gdef\csname CGTip@subsec:clock_intro\endcsname{Section 1.1: The Riemann--von Mangoldt clock and the ordinates (p. 3)}
\expandafter\gdef\csname CGTip@eq:intro_rvm\endcsname{Eq. (1.1), p. 3: N(T)=N\_0(T)+(1)/(pi)arg zeta ((1)/(2)+iT)+O((1)/(T)), N\_0(T)=(T)/(2pi)log(T)/(2pi e)+(7)/(8),}
\expandafter\gdef\csname CGTip@eq:intro_clock\endcsname{Eq. (1.2), p. 3: N\_0(tau\_k)=k-(1)/(2), tau\_k=(2pi (k-(11)/(8)))/(W\_0((k-(11)/(8))/e))(k>=1).}
\expandafter\gdef\csname CGTip@eq:intro_SX\endcsname{Eq. (1.3), p. 4: S\_X(t)=-(1)/(pi)sum\_p<= Xarg (1-p\textasciicircum{}-1/2-it), F\_X(t)=N\_0(t)+S\_X(t),}
\expandafter\gdef\csname CGTip@eq:intro_defect\endcsname{Eq. (1.4), p. 4: s\_k=k-(1)/(2)-N\_0(gamma\_k), u\_k(X)=s\_k-S\_X(gamma\_k)=k-(1)/(2)-F\_X(gamma\_k).}
\expandafter\gdef\csname CGTip@subsec:program\endcsname{Section 1.2: The approximation of the zeros by the clock and the primes (p. 4)}
\expandafter\gdef\csname CGTip@subsec:known_intro\endcsname{Section 1.3: Background (p. 5)}
\expandafter\gdef\csname CGTip@subsec:comparison\endcsname{Section 1.4: Comparison with the cited antecedents (p. 6)}
\expandafter\gdef\csname CGTip@tab:antecedents\endcsname{Table 1 (p. 6): Cited antecedents, conclusions proved here, and the scope of the difference.}
\expandafter\gdef\csname CGTip@tab:evaluates\endcsname{Table 2 (p. 8): What each approximation of the zeros evaluates, and what is proved about its error.}
\expandafter\gdef\csname CGTip@subsec:main_results\endcsname{Section 1.5: What this paper adds (p. 9)}
\expandafter\gdef\csname CGTip@eq:intro_main\endcsname{Eq. (1.5), p. 10: (1)/(no.I\_T)sum\_kin I\_Tu\_k(X)\textasciicircum{}2=(1)/(2pi\textasciicircum{}2)log(1)/(c)+O(1+sqrt(log(1/c))),}
\expandafter\gdef\csname CGTip@eq:intro_D\endcsname{Eq. (1.6), p. 10: (1)/(no.I\_T)sum\_kin I\_T(tau\_k\textasciicircum{}X-gamma\_k)\textasciicircum{}2=(log log log T+O(sqrt(log log log T)))/(log\textasciicircum{}2T).}
\expandafter\gdef\csname CGTip@eq:intro_Dprime\endcsname{Eq. (1.7), p. 11: (1)/(no.I\_T)sum\_kin I\_T|hat x\_k-gamma\_k|\textasciicircum{}4<<((1+log(1/c))\textasciicircum{}2)/(log\textasciicircum{}4T),}
\expandafter\gdef\csname CGTip@eq:intro_Dpp\endcsname{Eq. (1.8), p. 11: (1)/(no.I\_T)sum\_kin I\_T(tau\_k\textasciicircum{}T\textasciicircum{}c-gamma\_k)\textasciicircum{}2<<(1+log log log T)/(log\textasciicircum{}2T),< t\_k(T\textasciicircum{}c)\textasciicircum{}2>\_I\_T<<1+log log log T,}
\expandafter\gdef\csname CGTip@eq:intro_E\endcsname{Eq. (1.9), p. 11: limsup\_k->inf|tau\_k-gamma\_k|<=(1)/(log 2)int\_0\textasciicircum{}1(e(sigma) dsigma)/(sqrt(sigma(2-sigma))),}
\expandafter\gdef\csname CGTip@eq:intro_F\endcsname{Eq. (1.10), p. 12: (N\_0\textasciicircum{}s(T,2T))/(no.I\_T)>=(4)/(3)-4Q\_2(T,X)-o(1), Q\_2(T,X)=(1)/(no.I\_T)int\_T\textasciicircum{}2T(N-F\_X)\textasciicircum{}2dF\_X .}
\expandafter\gdef\csname CGTip@eq:intro_G\endcsname{Eq. (1.11), p. 12: zeta\_kappa(s)=zeta\_tau(s) (P\_X(s))/(P\_X(1)) e\textasciicircum{}H\_X(s), P\_X(s)=prod\_p<= X(1-p\textasciicircum{}-s)\textasciicircum{}-1,}
\expandafter\gdef\csname CGTip@eq:intro_ratio\endcsname{Eq. (1.12), p. 12: (<(tau\_k\textasciicircum{}X-gamma\_k)\textasciicircum{}2>\_I\_T)/(<(tau\_k-gamma\_k)\textasciicircum{}2>\_I\_T)\textasciitilde{}(log log log T)/(2log log T)->0,}
\expandafter\gdef\csname CGTip@subsec:status\endcsname{Section 1.6: Degree of proof (p. 13)}
\expandafter\gdef\csname CGTip@tab:status\endcsname{Table 3 (p. 13): Main results: degree of proof, role, and whether zeta or its zeros are evaluated.}
\expandafter\gdef\csname CGTip@sec:notation\endcsname{Section 2: Notation (p. 14)}
\expandafter\gdef\csname CGTip@subsec:symbols\endcsname{Section 2.1: List of symbols (p. 14)}
\expandafter\gdef\csname CGTip@subsec:counting\endcsname{Section 2.2: Counting functions (p. 16)}
\expandafter\gdef\csname CGTip@eq:RvM\endcsname{Eq. (2.1), p. 16: N(t)=(1)/(pi)vartheta(t)+1+S(t), S(t)=(1)/(pi)arg zeta ((1)/(2)+it)=O(log t),}
\expandafter\gdef\csname CGTip@eq:N0_def\endcsname{Eq. (2.2), p. 16: Theta\_0(t)=(t)/(2)log(t)/(2pi e)-(pi)/(8), N\_0(t)=(1)/(pi)Theta\_0(t)+1, varpi(t)=vartheta(t)-Theta\_0(t).}
\expandafter\gdef\csname CGTip@eq:N_decomposition\endcsname{Eq. (2.3), p. 16: N(t)=N\_0(t)+S(t)+(varpi(t))/(pi), L(t)=N\_0'(t)=(1)/(2pi)log(t)/(2pi).}
\expandafter\gdef\csname CGTip@eq:IT_size\endcsname{Eq. (2.4), p. 16: no.I\_T=TL(T)+O(T).}
\expandafter\gdef\csname CGTip@subsec:clock\endcsname{Section 2.3: The RvM clock and the displacements (p. 17)}
\expandafter\gdef\csname CGTip@eq:clock\endcsname{Eq. (2.5), p. 17: tau\_x=(2pi (x-(11)/(8)))/(W\_0((x-(11)/(8))/e)), (dtau\_x)/(dx)=(1)/(L(tau\_x))=(2pi)/(log(tau\_x/2pi)),}
\expandafter\gdef\csname CGTip@def:displacements\endcsname{Definition 2.1 (p. 17): For k>=1, s\_k=k-(1)/(2)-N\_0(gamma\_k) and delta\_k=tau\_k-gamma\_k.}
\expandafter\gdef\csname CGTip@eq:s_delta\endcsname{Eq. (2.6), p. 17: s\_k=L(xi\_k) delta\_kfor some xi\_k between gamma\_k and tau\_k .}
\expandafter\gdef\csname CGTip@lem:s_bounds\endcsname{Lemma 2.2 (p. 17): For every k>=1, N(gamma\_k\textasciicircum{}-)-N\_0(gamma\_k)+(1)/(2)<= s\_k<= N(gamma\_k\textasciicircum{}+)-N\_0(gamma\_k)-(1)/(2), and |s\_k-(varpi(gamma\_k))/(pi)|<=max (|S(gamma\_k\textasciicircum{}-)|,|S(gamma\_k\textasciicircum{}+)|)-(1)/(2) .}
\expandafter\gdef\csname CGTip@eq:s_bounds\endcsname{Eq. (2.7), p. 17: N(gamma\_k\textasciicircum{}-)-N\_0(gamma\_k)+(1)/(2)<= s\_k<= N(gamma\_k\textasciicircum{}+)-N\_0(gamma\_k)-(1)/(2),}
\expandafter\gdef\csname CGTip@eq:s_S_bound\endcsname{Eq. (2.8), p. 17: |s\_k-(varpi(gamma\_k))/(pi)|<=max (|S(gamma\_k\textasciicircum{}-)|,|S(gamma\_k\textasciicircum{}+)|)-(1)/(2) .}
\expandafter\gdef\csname CGTip@subsec:euler_phase\endcsname{Section 2.4: Euler phases and the corrected clock (p. 17)}
\expandafter\gdef\csname CGTip@eq:SA_def\endcsname{Eq. (2.9), p. 17: S\_A(t)=sum\_nin Aa\_nsin(tlog n), v\_A=(1)/(2)sum\_nin Aa\_n\textasciicircum{}2 ;}
\expandafter\gdef\csname CGTip@eq:SX_def\endcsname{Eq. (2.10), p. 17: S\_X(t)=-(1)/(pi)sum\_p<= Xarg (1-p\textasciicircum{}-1/2-it)=-(1)/(pi)sum\_p<= Xsum\_m>=1(sin(mtlog p))/(m p\textasciicircum{}m/2)=S\_A\_X(t).}
\expandafter\gdef\csname CGTip@eq:sigma_M_def\endcsname{Eq. (2.11), p. 18: |S\_X(t)|<=sigma\_X:=(1)/(pi)sum\_p<= Xsum\_m>=1(1)/(m p\textasciicircum{}m/2), |S\_X'(t)|<= M\_X:=(1)/(pi)sum\_p<= X(log p)/(sqrt(p)-1),}
\expandafter\gdef\csname CGTip@def:defect\endcsname{Definition 2.3 (p. 18): For X>=2 let F\_X=N\_0+S\_X and u\_k(X)=k-(1)/(2)-F\_X(gamma\_k)=s\_k-S\_X(gamma\_k). The one-step clock and its residual are tau\_k\textasciicircum{}X=tau\_k-(S\_X(tau\_k))/(L(tau\_k)), delta\_k\textasciicircum{}X=tau\_k\textasciicircum{}X-gamma\_k, r\_k(X)=L(tau\_k) delta\_k\textasciicircum{}X .}
\expandafter\gdef\csname CGTip@eq:one_step_clock\endcsname{Eq. (2.12), p. 18: tau\_k\textasciicircum{}X=tau\_k-(S\_X(tau\_k))/(L(tau\_k)), delta\_k\textasciicircum{}X=tau\_k\textasciicircum{}X-gamma\_k, r\_k(X)=L(tau\_k) delta\_k\textasciicircum{}X .}
\expandafter\gdef\csname CGTip@sec:known_results\endcsname{Section 3: External results (p. 18)}
\expandafter\gdef\csname CGTip@thm:known_basic\endcsname{Theorem 3.1 ([ref, Thms. 9.4, 9.9 and S 9.12], [ref]) (p. 18): S(t)=O(log t); gamma\_k+1-gamma\_k is bounded; and S\_1(t)=int\_0\textasciicircum{}tS(y) dy=O(log t).}
\expandafter\gdef\csname CGTip@thm:known_BW\endcsname{Theorem 3.2 (Bellotti--Wong [ref, Thm. 1.1, Cor. 1.5 and (1.2)]) (p. 18): For T>= e, |N(T)-(T)/(2pi)log(T)/(2pi e)|<=BW(T) and |S(T)|<=BWS(T), where BW(t) =min 0.10076log t+0.24460log log t+8.08344,; 0.11200log t+0.12567log log t+3.77417,; BWS(t) =0.10076log t+min 0.24460log log t+7.20844, 1.68845log log t+1.50956. Moreover |S(T)|<=2.5167 for 0<= T<=30 610 046 000.}
\expandafter\gdef\csname CGTip@eq:BWS_def\endcsname{Eq. (3.1), p. 18: BW(t) =min 0.10076log t+0.24460log log t+8.08344,; 0.11200log t+0.12567log log t+3.77417,; BWS(t) =0.10076log t+min 0.24460log log t+7.20844, 1.68845log log t+1.50956.}
\expandafter\gdef\csname CGTip@thm:known_selberg_moments\endcsname{Theorem 3.3 (Selberg [ref]; see [ref, p. 225]) (p. 18): For every integer j>=1, int\_T\textasciicircum{}2TS(t)\textasciicircum{}2j dt=kappa\_j T(log log T)\textasciicircum{}j+O\_j(T(log log T)\textasciicircum{}j-1/2).}
\expandafter\gdef\csname CGTip@thm:known_korolev\endcsname{Theorem 3.4 (Selberg; Korolev [ref], [ref, Thm. 4.2, p. 73]) (p. 19): Let Delta\_k be the integer such that g\_k+Delta\_k-2<gamma\_k<= g\_k+Delta\_k-1 (in [ref] the Gram points are t\_m=g\_m-1 and t\_m-1<gamma\_k<= t\_m with m=k+Delta\_k), and let 0<varepsilon\_0<10\textasciicircum{}-3. For every fixed j>=1, as n\_0->inf, uniformly for n\_0\textasciicircum{}27/82+0.9varepsilon\_0<= h\_0<= n\_0\textasciicircum{}27/82+varepsilon\_0, sum\_n\_0<k<= n\_0+h\_0Del...}
\expandafter\gdef\csname CGTip@eq:korolev_moments\endcsname{Eq. (3.2), p. 19: sum\_n\_0<k<= n\_0+h\_0Delta\_k\textasciicircum{}2j=kappa\_jh\_0(log log n\_0)\textasciicircum{}j(1+o(1)), |sum\_n\_0<k<= n\_0+h\_0Delta\_k\textasciicircum{}2j-1|<<\_j h\_0(log log n\_0)\textasciicircum{}j-1.}
\expandafter\gdef\csname CGTip@thm:known_fujii_selberg\endcsname{Theorem 3.5 (Selberg's formula with Fujii's bounds [ref, S 3, pp. 232--233]) (p. 19): Let a\_2 be fixed in the range of Fujii's Main Theorem [ref, p. 225], which contains (1)/(5). There are T\_0 and C, depending only on a\_2, such that for T>= T\_0, T\textasciicircum{}1/2+a\_2<= H<= T and every integer k>=1 with x\_rm F=T\textasciicircum{}a\_2/(60k)>=2, int\_T\textasciicircum{}T+H(S(t)-S\textasciicircum{}rm pr\_x\_rm F\textasciicircum{}3(t))\textasciicircum{}2kdt<=(Ck\textasciicircum{}2)\textasciicircum{}2kH .}
\expandafter\gdef\csname CGTip@eq:fujii_selberg\endcsname{Eq. (3.3), p. 19: int\_T\textasciicircum{}T+H(S(t)-S\textasciicircum{}rm pr\_x\_rm F\textasciicircum{}3(t))\textasciicircum{}2kdt<=(Ck\textasciicircum{}2)\textasciicircum{}2kH .}
\expandafter\gdef\csname CGTip@lem:selberg_uniform\endcsname{Lemma 3.6 (Selberg's approximation uniformly in the order) (p. 19): Let kappa\_rm F=1/(300log 2) and Y\_k=T\textasciicircum{}1/(80k). There are absolute constants C and T\_0 such that for T>= T\_0 and every integer k with 1<= k<=kappa\_rm Flog T, int\_T\textasciicircum{}2T(S(t)-S\textasciicircum{}rm pr\_Y\_k(t))\textasciicircum{}2kdt<=(Ck\textasciicircum{}2)\textasciicircum{}2k T . In particular int\_T\textasciicircum{}2T(S-S\textasciicircum{}rm pr\_Y\_0)\textasciicircum{}4<< T with Y\_0=Y\_2=T\textasciicircum{}1/160.}
\expandafter\gdef\csname CGTip@eq:selberg_uniform\endcsname{Eq. (3.4), p. 19: int\_T\textasciicircum{}2T(S(t)-S\textasciicircum{}rm pr\_Y\_k(t))\textasciicircum{}2kdt<=(Ck\textasciicircum{}2)\textasciicircum{}2k T .}
\expandafter\gdef\csname CGTip@rem:korolev_uniform\endcsname{Remark 3.7 (An independent route to the uniformity in the order) (p. 20): Korolev [ref, Lemma 3.13, p. S126] proves, under the conditions of his Lemma 3.11 (p. S124) and the conventions of pp. S13 and S123: for fixed 0<varepsilon\_rm K<10\textasciicircum{}-3, c\_0=2880varepsilon\_rm K\textasciicircum{}-1, T\_1>= T\_0(varepsilon\_rm K), T\_1\textasciicircum{}27/82+0.9varepsilon\_rm K<= H<= T\_1\textasciicircum{}27/82+varepsilon\_rm K, Y\_rm K'=T\_1\textasciicircum{}0.1varepsilon\_rm K,...}
\expandafter\gdef\csname CGTip@eq:korolev_lemma\endcsname{Eq. (3.5), p. 20: int\_T\_1\textasciicircum{}T\_1+H|S(t)+(1)/(pi)sum\_p<= Y\_rm K(sin(tlog p))/(sqrt(p))|\textasciicircum{}2kdt<(c\_0)/(56) (c\_0k)\textasciicircum{}2kH .}
\expandafter\gdef\csname CGTip@thm:known_fujii_difference\endcsname{Theorem 3.8 (Fujii [ref, (alpha'), p. 33]) (p. 20): There are absolute constants A and T\_0 such that for every integer k>=1, T>= T\_0 and 0<h<=1, [ int\_T\textasciicircum{}2T(S(t+h)-S(t))\textasciicircum{}2kdt<<(Ak)\textasciicircum{}kT(log ((hlog T+3) e\textasciicircum{}k))\textasciicircum{}k, ] with an absolute implied constant; the interval (T,2T) may be replaced by (0,T) [ref, p. 34, footnote].}
\expandafter\gdef\csname CGTip@eq:fujii_difference\endcsname{Eq. (3.6), p. 21: int\_T-1\textasciicircum{}2T(S(t+h)-S(t))\textasciicircum{}2kdt<= C T(log(2+hlog T))\textasciicircum{}k(k=1,2, 0<h<=1, T>= T\_0),}
\expandafter\gdef\csname CGTip@thm:known_MV\endcsname{Theorem 3.9 (Montgomery--Vaughan [ref, Cor. 3]) (p. 21): For complex z\_n, real t\_0 and H>0, int\_t\_0\textasciicircum{}t\_0+H|sum\_n<= Yz\_nn\textasciicircum{}-it|\textasciicircum{}2dt=sum\_n<= Y|z\_n|\textasciicircum{}2(H+O(n)) with an absolute implied constant.}
\expandafter\gdef\csname CGTip@thm:known_mertens\endcsname{Theorem 3.10 (Mertens [ref]; see [ref, Ch. I.1, Thms. 1.8 and 1.10]) (p. 21): For y>=2, sum\_p<= y(log p)/(p)=log y+O(1) and sum\_p<= y(1)/(p)=log log y+O(1).}
\expandafter\gdef\csname CGTip@thm:known_selberg_clt\endcsname{Theorem 3.11 (Selberg [ref]; Tsang [ref]; see [ref, (1.1)]) (p. 21): For fixed real y\_1<y\_2, (1)/(T)meastin[T,2T]:y\_1<S(t)/ s\_T<= y\_2->(1)/(sqrt(2pi))int\_y\_1\textasciicircum{}y\_2e\textasciicircum{}-y\textasciicircum{}2/2dy, where s\_T=((1)/(2pi\textasciicircum{}2)log log T)\textasciicircum{}1/2.}
\expandafter\gdef\csname CGTip@thm:known_maximal\endcsname{Theorem 3.12 (Hardy--Littlewood [ref, Thm. I.1], as quoted in [ref, (1.2)]; Kinnunen [ref, Thm. 1.4]) (p. 21): Let Mf(x)=sup\_y>0(1)/(2y)int\_x-y\textasciicircum{}x+y|f| be the centred maximal function of fin L\textasciicircum{}1\_rm loc(R). For 1<p<inf and fin L\textasciicircum{}p(R), || Mf||\_p<= A\_p||f||\_p, where A\_p is bounded for p>=2. If f and f' belong to L\textasciicircum{}p(R), 1<p<inf, then Mf has a locally absolutely continuous representative, (Mf)'in L\textasciicircum{}p(R), and |(Mf)'|<= M(f') almos...}
\expandafter\gdef\csname CGTip@thm:known_carleson\endcsname{Theorem 3.13 (Carleson [ref], Hunt [ref]; on R see [ref, (1.1) and p. 1]) (p. 21): For fin L\textasciicircum{}2(R)n L\textasciicircum{}4(R) with Fourier transform Ff(omega)=int\_R f(t)e\textasciicircum{}-iomega tdt, let Cf(x)=sup\_omega\_1in R|(1)/(2pi)int\_-inf\textasciicircum{}omega\_1 Ff(omega)e\textasciicircum{}iomega xdomega |. Then || Cf||\_4<<||f||\_4.}
\expandafter\gdef\csname CGTip@thm:known_exponents\endcsname{Theorem 3.14 ([ref, Lemma 2.8 and (7.57)]; Bourgain [ref, S 5, Thm. 6]; Heath-Brown [ref, Thm. 2]) (p. 21): (a) If (k, l) is an exponent pair, so is A(k, l)=((k)/(2 k+2),(1)/(2)+(l)/(2 k+2)). (b) If (k, l) is an exponent pair with l- k>=(1)/(2), then zeta(l- k+it)<<|t|\textasciicircum{} klog|t| for |t|>=2. (c) For every varepsilon>0, ((13)/(84)+varepsilon,(55)/(84)+varepsilon) is an exponent pair. (d) For every integer m>=3 and every vare...}
\expandafter\gdef\csname CGTip@thm:known_PL\endcsname{Theorem 3.15 (Phragmen--Lindelof principle in a strip [ref, Lemma 2.5], [ref, Lemma 3]) (p. 22): Let sigma\_1<sigma\_2 and e\_1>= e\_2>=0. If |zeta(sigma\_i+it)|<= C\_i|t|\textasciicircum{}e\_ilog|t| for |t|>=2 and i=1,2, then |zeta(sigma+it)|<<|t|\textasciicircum{}e(sigma)log|t| uniformly for sigma\_1<=sigma <=sigma\_2 and |t|>=2, where e is the linear function with e(sigma\_i)=e\_i.}
\expandafter\gdef\csname CGTip@thm:known_backlund\endcsname{Theorem 3.16 (Backlund's method [ref, Lemmas 2.1, 2.3 and 2.4], after [ref]) (p. 22): Let T be large, varkappa in N, f\_varkappa(s)=(1)/(2)[((s+iT-1)zeta(s+iT))\textasciicircum{}varkappa+((s-iT-1)zeta(s-iT))\textasciicircum{}varkappa ], and for a disc D(sigma\_o,varrho) let Z\_varkappa(sigma\_o,varrho)=(1)/(varkappa)sum\_zlog(varrho)/(|z-sigma\_o|) over the zeros z of f\_varkappa in D(sigma\_o,varrho). Let 1<sigma\_o<varrho, sigma\textasciicircum{}sharp=sigma...}
\expandafter\gdef\csname CGTip@rem:korolev\endcsname{Remark 3.17 (Known consequences for the clock) (p. 22): (a) Moments and law of s\_k. Since vartheta(gamma\_k)/pi in(k+Delta\_k-2,k+Delta\_k-1] and vartheta/pi=N\_0-1+varpi/pi, one has s\_k+Delta\_k=k+Delta\_k-(1)/(2)-N\_0(gamma\_k)in[-(1)/(2)+(varpi(gamma\_k))/(pi),(1)/(2)+(varpi(gamma\_k))/(pi)), so s\_k=-Delta\_k+ d\_k with | d\_k|<1. The index set I\_T is the integer interval (N(T),N(...}
\expandafter\gdef\csname CGTip@eq:s_moments\endcsname{Eq. (3.7), p. 22: < s\_k\textasciicircum{}2j>\_I\_T=kappa\_j(log log T)\textasciicircum{}j(1+o(1)), < s\_k\textasciicircum{}2j-1>\_I\_T=O\_j((log log T)\textasciicircum{}j-1).}
\expandafter\gdef\csname CGTip@eq:BW_pointwise\endcsname{Eq. (3.8), p. 23: -(BW(gamma\_k)+3/8)/(L(tau\_k))<=delta\_k<=(BW(gamma\_k)-11/8)/(L(gamma\_k))(k>=1).}
\expandafter\gdef\csname CGTip@sec:identities\endcsname{Section 4: Theorem A: an exact defect identity at the zeros (p. 23)}
\expandafter\gdef\csname CGTip@thm:identity_general\endcsname{Theorem 4.1 (p. 23): Let 1<= k\_1<k\_2 be indices, F a real C\textasciicircum{}1 function on an open interval containing [gamma\_k\_1,gamma\_k\_2], and u\_k=k-(1)/(2)-F(gamma\_k) for k\_1<= k<= k\_2. For every integer j>=1, int\_gamma\_k\_1\textasciicircum{}gamma\_k\_2(N(t)-F(t))\textasciicircum{}2jF'(t) dt=sum\_k=k\_1+1\textasciicircum{}k\_2-1Psi\_j(u\_k)+(((1)/(2)+u\_k\_1)\textasciicircum{}2j+1+((1)/(2)-u\_k\_2)\textasciicircum{}2j+1)/(2j+1), where Psi\_j(u)=...}
\expandafter\gdef\csname CGTip@eq:identity_general\endcsname{Eq. (4.1), p. 23: int\_gamma\_k\_1\textasciicircum{}gamma\_k\_2(N(t)-F(t))\textasciicircum{}2jF'(t) dt=sum\_k=k\_1+1\textasciicircum{}k\_2-1Psi\_j(u\_k)+(((1)/(2)+u\_k\_1)\textasciicircum{}2j+1+((1)/(2)-u\_k\_2)\textasciicircum{}2j+1)/(2j+1),}
\expandafter\gdef\csname CGTip@eq:Psi_def\endcsname{Eq. (4.2), p. 23: Psi\_j(u)=(((1)/(2)+u)\textasciicircum{}2j+1+((1)/(2)-u)\textasciicircum{}2j+1)/(2j+1)=sum\_l=0\textasciicircum{}jbinom(2j+1,2l)(4\textasciicircum{}l-j)/(2j+1) u\textasciicircum{}2l=u\textasciicircum{}2j+sum\_l<jbinom(2j+1,2l)(4\textasciicircum{}l-j)/(2j+1)u\textasciicircum{}2l.}
\expandafter\gdef\csname CGTip@eq:identity_quadratic\endcsname{Eq. (4.3), p. 23: sum\_k=k\_1+1\textasciicircum{}k\_2-1u\_k\textasciicircum{}2=int\_gamma\_k\_1\textasciicircum{}gamma\_k\_2(N(t)-F(t))\textasciicircum{}2F'(t) dt-(k\_2-k\_1-1)/(12)-(((1)/(2)+u\_k\_1)\textasciicircum{}3+((1)/(2)-u\_k\_2)\textasciicircum{}3)/(3) .}
\expandafter\gdef\csname CGTip@eq:identity_one_gap\endcsname{Eq. (4.4), p. 23: int\_gamma\_k\textasciicircum{}gamma\_k+1(k-F(t))\textasciicircum{}2jF'(t) dt =((k-F(gamma\_k))\textasciicircum{}2j+1-(k-F(gamma\_k+1))\textasciicircum{}2j+1)/(2j+1); =(((1)/(2)+u\_k)\textasciicircum{}2j+1+((1)/(2)-u\_k+1)\textasciicircum{}2j+1)/(2j+1),}
\expandafter\gdef\csname CGTip@thm:identity_g\endcsname{Theorem 4.2 (p. 24): Let k\_1<k\_2, F and u\_k be as in Theorem 4.1, and let g:R -> C be continuous with a primitive G. Then int\_gamma\_k\_1\textasciicircum{}gamma\_k\_2 g(N(t)-F(t))F'(t) dt=sum\_k=k\_1+1\textasciicircum{}k\_2-1[ G(u\_k+(1)/(2))- G(u\_k-(1)/(2))]+ G(u\_k\_1+(1)/(2))- G(u\_k\_2-(1)/(2)).}
\expandafter\gdef\csname CGTip@eq:identity_g\endcsname{Eq. (4.5), p. 24: int\_gamma\_k\_1\textasciicircum{}gamma\_k\_2 g(N(t)-F(t))F'(t) dt=sum\_k=k\_1+1\textasciicircum{}k\_2-1[ G(u\_k+(1)/(2))- G(u\_k-(1)/(2))]+ G(u\_k\_1+(1)/(2))- G(u\_k\_2-(1)/(2)).}
\expandafter\gdef\csname CGTip@cor:identity_exp\endcsname{Corollary 4.3 (p. 24): For real omega !=0, (2sin(omega/2))/(omega)sum\_k=k\_1+1\textasciicircum{}k\_2-1e\textasciicircum{}iomega u\_k=int\_gamma\_k\_1\textasciicircum{}gamma\_k\_2e\textasciicircum{}iomega(N(t)-F(t))F'(t) dt-(e\textasciicircum{}iomega(u\_k\_1+1/2)-e\textasciicircum{}iomega(u\_k\_2-1/2))/(iomega).}
\expandafter\gdef\csname CGTip@eq:identity_exp\endcsname{Eq. (4.6), p. 24: (2sin(omega/2))/(omega)sum\_k=k\_1+1\textasciicircum{}k\_2-1e\textasciicircum{}iomega u\_k=int\_gamma\_k\_1\textasciicircum{}gamma\_k\_2e\textasciicircum{}iomega(N(t)-F(t))F'(t) dt-(e\textasciicircum{}iomega(u\_k\_1+1/2)-e\textasciicircum{}iomega(u\_k\_2-1/2))/(iomega).}
\expandafter\gdef\csname CGTip@cor:signed_law\endcsname{Corollary 4.4 (Signed law of s\_k from the t-law of S) (p. 24): For fixed real xi, < e\textasciicircum{}ixi s\_k/ s\_T>\_I\_T-> e\textasciicircum{}-xi\textasciicircum{}2/2. Hence the empirical law of s\_k/ s\_T on I\_T tends to the standard normal law.}
\expandafter\gdef\csname CGTip@rem:twelfth\endcsname{Remark 4.5 (p. 24): Psi\_j(u)=int\_u-1/2\textasciicircum{}u+1/2x\textasciicircum{}2jdx is the mean of x\textasciicircum{}2j over the unit interval centred at u; in particular Psi\_1(u)=u\textasciicircum{}2+(1)/(12) and (1)/(12)=int\_-1/2\textasciicircum{}1/2x\textasciicircum{}2dx. The constant comes from the unit jump of N-F at each zero, through the primitive in Theorem 4.2; when F is not monotone, N-F need not run through an interval of...}
\expandafter\gdef\csname CGTip@cor:identity_clock\endcsname{Corollary 4.6 (p. 25): For X>=2 and 1<= k\_1<k\_2, sum\_k=k\_1+1\textasciicircum{}k\_2-1u\_k(X)\textasciicircum{}2= int\_gamma\_k\_1\textasciicircum{}gamma\_k\_2(S-S\_X+(varpi)/(pi))\textasciicircum{}2(L+S\_X')dt; -(k\_2-k\_1-1)/(12)-(((1)/(2)+u\_k\_1(X))\textasciicircum{}3+((1)/(2)-u\_k\_2(X))\textasciicircum{}3)/(3) .}
\expandafter\gdef\csname CGTip@eq:identity_corrected\endcsname{Eq. (4.7), p. 25: sum\_k=k\_1+1\textasciicircum{}k\_2-1u\_k(X)\textasciicircum{}2= int\_gamma\_k\_1\textasciicircum{}gamma\_k\_2(S-S\_X+(varpi)/(pi))\textasciicircum{}2(L+S\_X')dt; -(k\_2-k\_1-1)/(12)-(((1)/(2)+u\_k\_1(X))\textasciicircum{}3+((1)/(2)-u\_k\_2(X))\textasciicircum{}3)/(3) .}
\expandafter\gdef\csname CGTip@lem:first_moment_identity\endcsname{Lemma 4.7 (First-moment identity) (p. 25): Let 2pi<t\_1<t\_2 be points which are not ordinates and psi in C\textasciicircum{}1[t\_1,t\_2], real or complex. Then sum\_t\_1<gamma\_k<= t\_2s\_kpsi(gamma\_k)=int\_t\_1\textasciicircum{}t\_2G(t)psi(t)L(t) dt+(1)/(2)[psi G\textasciicircum{}2]\_t\_1\textasciicircum{}t\_2-(1)/(2)int\_t\_1\textasciicircum{}t\_2G(t)\textasciicircum{}2psi'(t) dt .}
\expandafter\gdef\csname CGTip@eq:first_moment_identity\endcsname{Eq. (4.8), p. 25: sum\_t\_1<gamma\_k<= t\_2s\_kpsi(gamma\_k)=int\_t\_1\textasciicircum{}t\_2G(t)psi(t)L(t) dt+(1)/(2)[psi G\textasciicircum{}2]\_t\_1\textasciicircum{}t\_2-(1)/(2)int\_t\_1\textasciicircum{}t\_2G(t)\textasciicircum{}2psi'(t) dt .}
\expandafter\gdef\csname CGTip@eq:first_moment_measures\endcsname{Eq. (4.9), p. 25: d(G\textasciicircum{}2)=-2GL dt+sum\_t\_1<gamma<t\_22G\textasciicircum{}*(gamma)nu(gamma) delta\_gamma ,}
\expandafter\gdef\csname CGTip@sec:main_theorem\endcsname{Section 5: Theorem B: how much prime information the clock needs (p. 25)}
\expandafter\gdef\csname CGTip@thm:main\endcsname{Theorem 5.1 (p. 25): There are absolute constants T\_0 and C such that for T>= T\_0 and every X with 2<= X<= T\textasciicircum{}1/2, writing c=log X/log T, |(1)/(no.I\_T)sum\_kin I\_Tu\_k(X)\textasciicircum{}2-(1)/(2pi\textasciicircum{}2)log(1)/(c)|<= C(1+sqrt(log(1/c))).}
\expandafter\gdef\csname CGTip@eq:main_theorem\endcsname{Eq. (5.1), p. 25: |(1)/(no.I\_T)sum\_kin I\_Tu\_k(X)\textasciicircum{}2-(1)/(2pi\textasciicircum{}2)log(1)/(c)|<= C(1+sqrt(log(1/c))).}
\expandafter\gdef\csname CGTip@cor:fixed_cutoff\endcsname{Corollary 5.2 (fixed cutoff) (p. 25): For fixed X>=2, (1)/(no.I\_T)sum\_kin I\_Tu\_k(X)\textasciicircum{}2=(1)/(2pi\textasciicircum{}2)log log T+O\_X(sqrt(log log T)), and no.kin I\_T:|u\_k(X)|<= y\_0=o(no.I\_T) for every fixed y\_0>0.}
\expandafter\gdef\csname CGTip@cor:power_cutoff\endcsname{Corollary 5.3 (power cutoff) (p. 26): For fixed 0<c<=(1)/(2) and X=T\textasciicircum{}c, limsup\_T->inf(1)/(no.I\_T)sum\_kin I\_Tu\_k(T\textasciicircum{}c)\textasciicircum{}2<=(1)/(2pi\textasciicircum{}2)log(1)/(c)+C(1+sqrt(log(1/c))), and no.kin I\_T:|u\_k(T\textasciicircum{}c)|>w\_T=o(no.I\_T) whenever w\_T->inf. The same holds for X=N(T)\textasciicircum{}c with fixed 0<c<(1)/(2), since then X=T\textasciicircum{}c' with c'=c+O(log log T/log T).}
\expandafter\gdef\csname CGTip@eq:power_cutoff\endcsname{Eq. (5.2), p. 26: limsup\_T->inf(1)/(no.I\_T)sum\_kin I\_Tu\_k(T\textasciicircum{}c)\textasciicircum{}2<=(1)/(2pi\textasciicircum{}2)log(1)/(c)+C(1+sqrt(log(1/c))),}
\expandafter\gdef\csname CGTip@cor:dichotomy\endcsname{Corollary 5.4 (dichotomy) (p. 26): Let X=X(T) with 2<= X<= T\textasciicircum{}1/2 for T>= T\_0. Then [ sup\_T>= T\_0(1)/(no.I\_T)sum\_kin I\_Tu\_k(X(T))\textasciicircum{}2<inf if and only ifliminf\_T->inf(log X)/(log T)>0 . ] For fixed X\_0>=2 and cin(0,(1)/(2)], (1)/(no.I\_T)sum\_kin I\_T(u\_k(X\_0)\textasciicircum{}2-u\_k(T\textasciicircum{}c)\textasciicircum{}2)=(1)/(2pi\textasciicircum{}2)log log T+O(sqrt(log log T)) .}
\expandafter\gdef\csname CGTip@eq:dichotomy\endcsname{Eq. (5.3), p. 26: (1)/(no.I\_T)sum\_kin I\_T(u\_k(X\_0)\textasciicircum{}2-u\_k(T\textasciicircum{}c)\textasciicircum{}2)=(1)/(2pi\textasciicircum{}2)log log T+O(sqrt(log log T)) .}
\expandafter\gdef\csname CGTip@prop:rounding\endcsname{Proposition 5.5 (Rounding the corrected count misassigns a positive proportion) (p. 26): There are absolute constants c\_*>0, kappa\_0>0 and T\_0 such that for T>= T\_0 and 2<= X<= T\textasciicircum{}c\_*, [ (1)/(T)meastin[T,2T]:floor(F\_X(t)+(1)/(2))!= N(t)>=kappa\_0 . ]}
\expandafter\gdef\csname CGTip@rem:comparison_cutoff\endcsname{Remark 5.6 (p. 26): Since logno.p<= X=log X+O(log log X), the condition of Corollary 5.4 says that the logarithm of the number of primes must grow proportionally to log T. A fixed number of primes leaves almost all zeros violating the corrected equation by more than any fixed amount; growth of the number of primes alone is necessary but...}
\expandafter\gdef\csname CGTip@rem:localization\endcsname{Remark 5.7 (Localization of the ordinates) (p. 27): If M\_X<=(1)/(2)L(T), then F\_X'>=L-M\_X>=(1)/(2)L(T) on [T,2T], so F\_X is strictly increasing there. For kin I\_T with F\_X(T)<= k-(1)/(2)<= F\_X(2T) the equation F\_X(t)=k-(1)/(2) has exactly one root t\_k\textasciicircum{}Xin[T,2T] (intermediate value theorem), and u\_k(X)=F\_X(t\_k\textasciicircum{}X)-F\_X(gamma\_k) gives |gamma\_k-t\_k\textasciicircum{}X|<=2|u\_k(X)|/L(T) (mea...}
\expandafter\gdef\csname CGTip@subsec:main_lemmas\endcsname{Section 5.1: Lemmas on prime polynomials (p. 27)}
\expandafter\gdef\csname CGTip@eq:D_def\endcsname{Eq. (5.4), p. 27: W\_1(t)=-(1)/(pi)sum\_pin J(sin(tlog p))/(sqrt(p)), V=sum\_pin J(1)/(p),; W=S\textasciicircum{}rm pr\_Y\_0-S\_X=+- W\_1-S\_X\textasciicircum{}(2)-S\_X\textasciicircum{}(>=3),}
\expandafter\gdef\csname CGTip@lem:V_size\endcsname{Lemma 5.8 (p. 27): V=log(1/c)+O(1) uniformly for log 2/log T<= c<=(1)/(2).}
\expandafter\gdef\csname CGTip@lem:trig_MV\endcsname{Lemma 5.9 (p. 27): For real z\_n, 2<= n<= Y, int\_T\textasciicircum{}2T(sum\_n<= Yz\_nsin(tlog n))\textasciicircum{}2dt=(T)/(2)sum\_n<= Yz\_n\textasciicircum{}2+O(sum\_n<= Ynz\_n\textasciicircum{}2+(sum\_n<= Y|z\_n|)\textasciicircum{}2), and the same with cos in place of sin.}
\expandafter\gdef\csname CGTip@eq:trig_MV\endcsname{Eq. (5.5), p. 27: int\_T\textasciicircum{}2T(sum\_n<= Yz\_nsin(tlog n))\textasciicircum{}2dt=(T)/(2)sum\_n<= Yz\_n\textasciicircum{}2+O(sum\_n<= Ynz\_n\textasciicircum{}2+(sum\_n<= Y|z\_n|)\textasciicircum{}2),}
\expandafter\gdef\csname CGTip@lem:D_moments\endcsname{Lemma 5.10 (p. 27): With absolute implied constants, int\_T\textasciicircum{}2T W\textasciicircum{}2 dt =(T)/(2pi\textasciicircum{}2)V+O(T(1+sqrt(V))), ; int\_T\textasciicircum{}2T W\textasciicircum{}4 dt << T(1+V)\textasciicircum{}2, ; int\_T\textasciicircum{}2TS\_X'(t)\textasciicircum{}2 dt << T(1+log X)\textasciicircum{}2 .}
\expandafter\gdef\csname CGTip@eq:D_second\endcsname{Eq. (5.6), p. 27: int\_T\textasciicircum{}2T W\textasciicircum{}2 dt =(T)/(2pi\textasciicircum{}2)V+O(T(1+sqrt(V))),}
\expandafter\gdef\csname CGTip@eq:D_fourth\endcsname{Eq. (5.7), p. 27: int\_T\textasciicircum{}2T W\textasciicircum{}4 dt << T(1+V)\textasciicircum{}2,}
\expandafter\gdef\csname CGTip@eq:SX_derivative\endcsname{Eq. (5.8), p. 27: int\_T\textasciicircum{}2TS\_X'(t)\textasciicircum{}2 dt << T(1+log X)\textasciicircum{}2 .}
\expandafter\gdef\csname CGTip@subsec:main_proof\endcsname{Section 5.2: Proof of Theorem 5.1 (p. 28)}
\expandafter\gdef\csname CGTip@eq:E0_bounds\endcsname{Eq. (5.9), p. 28: int\_T\textasciicircum{}2TE\_0\textasciicircum{}4 dt<< T,int\_T\textasciicircum{}2TE\_0\textasciicircum{}2 dt<= T\textasciicircum{}1/2(int\_T\textasciicircum{}2TE\_0\textasciicircum{}4 dt)\textasciicircum{}1/2<< T .}
\expandafter\gdef\csname CGTip@eq:J_asymptotic\endcsname{Eq. (5.10), p. 28: E\_2=(T)/(2pi\textasciicircum{}2)V+O(T(1+sqrt(V))),}
\expandafter\gdef\csname CGTip@eq:R_fourth\endcsname{Eq. (5.11), p. 28: int\_T\textasciicircum{}2TE\_X\textasciicircum{}4 dt<< T(1+V)\textasciicircum{}2 .}
\expandafter\gdef\csname CGTip@eq:main_step1\endcsname{Eq. (5.12), p. 28: sum\_kin I\_Tu\_k(X)\textasciicircum{}2=int\_T\textasciicircum{}2TE\_X\textasciicircum{}2L dt+int\_T\textasciicircum{}2TE\_X\textasciicircum{}2S\_X' dt-(no.I\_T)/(12)+O(T\textasciicircum{}3/4).}
\expandafter\gdef\csname CGTip@eq:main_step2\endcsname{Eq. (5.13), p. 29: |int\_T\textasciicircum{}2TE\_X\textasciicircum{}2S\_X' dt|<=(int\_T\textasciicircum{}2TE\_X\textasciicircum{}4)\textasciicircum{}1/2(int\_T\textasciicircum{}2TS\_X'\textasciicircum{}2)\textasciicircum{}1/2<< T(1+V)(1+log X).}
\expandafter\gdef\csname CGTip@eq:main_step3\endcsname{Eq. (5.14), p. 29: (1)/(no.I\_T)sum\_kin I\_Tu\_k(X)\textasciicircum{}2=(V)/(2pi\textasciicircum{}2)+O(1+sqrt(V)+(1+V)/(log T)+(1+V) c),}
\expandafter\gdef\csname CGTip@subsec:variable_cutoff\endcsname{Section 5.3: A cutoff varying with the index (p. 29)}
\expandafter\gdef\csname CGTip@lem:maximal_primes\endcsname{Lemma 5.11 (p. 29): Let T>= T\_0 and 2<= X\_-<X\_+<=(T/3)\textasciicircum{}1/2 with X\_+<=4X\_-; let V\_+-=sum\_X\_-<p<= X\_+p\textasciicircum{}-1 and [ Q(t)=max\_X\_-<= y<= X\_+|sum\_X\_-<p<= yp\textasciicircum{}-1/2-it| . ] Then Q is Lipschitz, | Q'(t)|<=max\_X\_-<= y<= X\_+|sum\_X\_-<p<= ylog p p\textasciicircum{}-1/2-it| for almost every t, and int\_T-3\textasciicircum{}2T+3 Q\textasciicircum{}4 dt<< TV\_+-\textasciicircum{}2,int\_T-3\textasciicircum{}2T+3| Q'|\textasciicircum{}4 dt<< T(log X\_+)\textasciicircum{}4V\_+-\textasciicircum{}2...}
\expandafter\gdef\csname CGTip@eq:maximal_L4\endcsname{Eq. (5.15), p. 29: int\_T-3\textasciicircum{}2T+3 Q\textasciicircum{}4 dt<< TV\_+-\textasciicircum{}2,int\_T-3\textasciicircum{}2T+3| Q'|\textasciicircum{}4 dt<< T(log X\_+)\textasciicircum{}4V\_+-\textasciicircum{}2,}
\expandafter\gdef\csname CGTip@eq:maximal_points\endcsname{Eq. (5.16), p. 29: < Q(gamma\_k)\textasciicircum{}2>\_I\_T<< V\_+-(1+(log X\_+)/(log T)sqrt(log log T))+(X\_+)/(T),< Q(tau\_k)\textasciicircum{}2>\_ I\textasciicircum{}tau\_T<< V\_+- .}
\expandafter\gdef\csname CGTip@eq:maximal_stieltjes\endcsname{Eq. (5.17), p. 30: sum\_kin I\_T Q(gamma\_k)\textasciicircum{}2=int\_T\textasciicircum{}2T Q\textasciicircum{}2L dt+[ Q\textasciicircum{}2G]\_T\textasciicircum{}2T-2int\_T\textasciicircum{}2TG Q Q' dt,}
\expandafter\gdef\csname CGTip@thm:variable_cutoff\endcsname{Theorem 5.12 (p. 30): For (log 2)/(log T)<= c<=(1)/(2) let X\_k=max 2,(tau\_k/2pi)\textasciicircum{}c and tilde u\_k=s\_k-S\_X\_k(gamma\_k). Uniformly in c, |<tilde u\_k\textasciicircum{}2>\_I\_T-(1)/(2pi\textasciicircum{}2)log(1)/(c)|<<1+sqrt(log(1/c)) . For fixed c, <(tilde u\_k-u\_k(X\_-))\textasciicircum{}2>\_I\_T->0, with X\_- as in the proof.}
\expandafter\gdef\csname CGTip@eq:variable_cutoff\endcsname{Eq. (5.18), p. 30: |<tilde u\_k\textasciicircum{}2>\_I\_T-(1)/(2pi\textasciicircum{}2)log(1)/(c)|<<1+sqrt(log(1/c)) .}
\expandafter\gdef\csname CGTip@sec:accounting\endcsname{Section 6: Theorem C: exact accounting for the one-step clock (p. 31)}
\expandafter\gdef\csname CGTip@eq:two_responses\endcsname{Eq. (6.1), p. 31: R\textasciicircum{}gamma\_n(I)=(2)/(a\_n)< s\_ksin(gamma\_klog n)>\_I, R\textasciicircum{}tau\_n(I)=(2)/(a\_n)< s\_ksin(tau\_klog n)>\_I,}
\expandafter\gdef\csname CGTip@eq:Cr_def\endcsname{Eq. (6.2), p. 31: Cr\_A(I)=< S\_A(tau\_k)\textasciicircum{}2>\_I-v\_A .}
\expandafter\gdef\csname CGTip@thm:accounting\endcsname{Theorem 6.1 (Finite-window accounting) (p. 31): Let I be a finite nonempty set of indices and A a set of prime powers with sum\_nin A|a\_n|<inf. item For every k and X>=2, r\_k(X)=u\_k(X)+(S\_X(gamma\_k)-S\_X(tau\_k))+chi\_k,; chi\_k=(gamma\_k)/(2pi)((tau\_k)/(gamma\_k)-1-log(tau\_k)/(gamma\_k)),(delta\_k\textasciicircum{}2)/(4pi max(gamma\_k,tau\_k))<=chi\_k<=(delta\_k\textasciicircum{}2)/(4pi min(gamma\_k,tau\_k)),;...}
\expandafter\gdef\csname CGTip@eq:one_step_decomposition\endcsname{Eq. (6.3), p. 31: r\_k(X)=u\_k(X)+(S\_X(gamma\_k)-S\_X(tau\_k))+chi\_k,; chi\_k=(gamma\_k)/(2pi)((tau\_k)/(gamma\_k)-1-log(tau\_k)/(gamma\_k)),(delta\_k\textasciicircum{}2)/(4pi max(gamma\_k,tau\_k))<=chi\_k<=(delta\_k\textasciicircum{}2)/(4pi min(gamma\_k,tau\_k)),; |S\_X(gamma\_k)-S\_X(tau\_k)|<= M\_X|delta\_k|.}
\expandafter\gdef\csname CGTip@eq:exact_gain\endcsname{Eq. (6.4), p. 31: < s\_k\textasciicircum{}2-(s\_k-lambda S\_A(tau\_k))\textasciicircum{}2>\_I =sum\_nin Aa\_n\textasciicircum{}2(lambda R\textasciicircum{}tau\_n(I)-(lambda\textasciicircum{}2)/(2))-lambda\textasciicircum{}2Cr\_A(I).}
\expandafter\gdef\csname CGTip@eq:full_factor_gain\endcsname{Eq. (6.5), p. 31: < s\_k\textasciicircum{}2-(s\_k-S\_X(tau\_k))\textasciicircum{}2>\_I=D\_X(I)-Cr\_X(I),; D\_X(I)=sum\_nin A\_Xa\_n\textasciicircum{}2(R\textasciicircum{}tau\_n(I)-(1)/(2))=2< s\_kS\_X(tau\_k)>\_I-v\_X,}
\expandafter\gdef\csname CGTip@eq:physical_gain\endcsname{Eq. (6.6), p. 32: <(L(tau\_k)delta\_k)\textasciicircum{}2-r\_k(X)\textasciicircum{}2>\_I=D\_X(I)-Cr\_X(I)+2<chi\_kS\_X(tau\_k)>\_I,; |<chi\_kS\_X(tau\_k)>\_I|<=sigma\_Xmax\_kin Ichi\_k .}
\expandafter\gdef\csname CGTip@rem:accounting_meaning\endcsname{Remark 6.2 (p. 32): (6.5) is an exact sign criterion for the gain of the full Euler phase at the clock points: the gain is positive exactly when the diagonal responses outweigh the cross term. No independence of the phases tau\_klog p is assumed. A threshold R\textasciicircum{}tau\_p=(1)/(2) describes a single diagonal prime contribution and does not det...}
\expandafter\gdef\csname CGTip@rem:matrix_projection\endcsname{Remark 6.3 (Optimal trigonometric correction) (p. 32): For a finite set A of frequencies let varsigma\_n(k)=sin(tau\_klog n), Q=(<varsigma\_n\_1varsigma\_n\_2>\_I)\_n\_1,n\_2in A, b=(< s\_kvarsigma\_n(k)>\_I)\_nin A and Upsilon\_lambda=sum\_nlambda\_nvarsigma\_n. Ordinary least squares gives <(s\_k-Upsilon\_lambda(k))\textasciicircum{}2>\_I=< s\_k\textasciicircum{}2>\_I-2lambda\textasciicircum{} T b+lambda\textasciicircum{} T Qlambda, minimized by lambda= Q\textasciicircum{}d...}
\expandafter\gdef\csname CGTip@rem:transport\endcsname{Remark 6.4 (Clock point against zero) (p. 32): For a prime power n, sin(gamma\_kl\_n)-sin(tau\_kl\_n)=int\_0\textasciicircum{}1(d)/(dy)sin((tau\_k-ydelta\_k)l\_n) dy, whence, exchanging the finite mean and the integral, R\textasciicircum{}gamma\_n(I)-R\textasciicircum{}tau\_n(I)=-(2l\_n)/(a\_n)int\_0\textasciicircum{}1< s\_kdelta\_kcos ((tau\_k-ydelta\_k)l\_n)>\_I dy . For fixed n and I=I\_T the right side is O(<|s\_kdelta\_k|>\_I\_T)=O(log log T/log T...}
\expandafter\gdef\csname CGTip@eq:response_transport\endcsname{Eq. (6.7), p. 32: R\textasciicircum{}gamma\_n(I)-R\textasciicircum{}tau\_n(I)=-(2l\_n)/(a\_n)int\_0\textasciicircum{}1< s\_kdelta\_kcos ((tau\_k-ydelta\_k)l\_n)>\_I dy .}
\expandafter\gdef\csname CGTip@rem:fixed_sup\endcsname{Remark 6.5 (Fixed prime sets) (p. 32): For fixed X, |delta\_k\textasciicircum{}X-delta\_k|<=sigma\_X/L(tau\_k)->0 and |u\_k(X)-s\_k|<=sigma\_X. A fixed set of primes changes neither limsup\_k|delta\_k| nor the unboundedness of the defect.}
\expandafter\gdef\csname CGTip@sec:computable\endcsname{Section 7: Theorem D: a computable clock that approaches the zeros (p. 33)}
\expandafter\gdef\csname CGTip@eq:r_decomposition\endcsname{Eq. (7.1), p. 33: r\_k(X)=u\_k(X)+ t\_k(X)+chi\_k, t\_k(X)=S\_X(gamma\_k)-S\_X(tau\_k),}
\expandafter\gdef\csname CGTip@lem:SX_fourth\endcsname{Lemma 7.1 (p. 33): For T>= T\_0 and 2<= X<= T\textasciicircum{}1/2, [ int\_T-3\textasciicircum{}2T+3S\_X'(t)\textasciicircum{}4 dt<< T(1+log X)\textasciicircum{}4,int\_T-3\textasciicircum{}2T+3S\_X''(t)\textasciicircum{}4 dt<< T(1+log X)\textasciicircum{}8 . ]}
\expandafter\gdef\csname CGTip@lem:transport\endcsname{Lemma 7.2 (Transport) (p. 33): For T>= T\_0 and 2<= X<= T\textasciicircum{}1/2, < t\_k(X)\textasciicircum{}2>\_I\_T<<(log log T (1+log X)\textasciicircum{}2)/(log\textasciicircum{}2T).}
\expandafter\gdef\csname CGTip@eq:transport\endcsname{Eq. (7.2), p. 33: < t\_k(X)\textasciicircum{}2>\_I\_T<<(log log T (1+log X)\textasciicircum{}2)/(log\textasciicircum{}2T).}
\expandafter\gdef\csname CGTip@eq:delta_fourth\endcsname{Eq. (7.3), p. 33: <delta\_k\textasciicircum{}4>\_I\_T=< s\_k\textasciicircum{}4L(xi\_k)\textasciicircum{}-4>\_I\_T<<((log log T)\textasciicircum{}2)/(log\textasciicircum{}4T).}
\expandafter\gdef\csname CGTip@eq:transport_maximal\endcsname{Eq. (7.4), p. 33: | t\_k(X)|=|int\_tau\_k\textasciicircum{}gamma\_k h(y) dy|<=int\_tau\_k-|delta\_k|\textasciicircum{}tau\_k+|delta\_k|| h|<=2|delta\_k| M h(tau\_k)}
\expandafter\gdef\csname CGTip@thm:computable_clock\endcsname{Theorem 7.3 (Theorem D) (p. 34): There are absolute constants T\_0 and C such that for T>= T\_0 and 2<= X<= T\textasciicircum{}1/2, with c=log X/log T, |< r\_k(X)\textasciicircum{}2>\_I\_T-(1)/(2pi\textasciicircum{}2)log(1)/(c)|<= C(1+sqrt(log(1/c)))(1+c\textasciicircum{}2log log T).}
\expandafter\gdef\csname CGTip@eq:computable_clock\endcsname{Eq. (7.5), p. 34: |< r\_k(X)\textasciicircum{}2>\_I\_T-(1)/(2pi\textasciicircum{}2)log(1)/(c)|<= C(1+sqrt(log(1/c)))(1+c\textasciicircum{}2log log T).}
\expandafter\gdef\csname CGTip@cor:computable_clock\endcsname{Corollary 7.4 (p. 34): Let X=T\textasciicircum{}1/sqrt(log log T). Then < r\_k(X)\textasciicircum{}2>\_I\_T =(1)/(4pi\textasciicircum{}2)log log log T+O(sqrt(log log log T)), ; (1)/(no.I\_T)sum\_kin I\_T(tau\_k\textasciicircum{}X-gamma\_k)\textasciicircum{}2 =(log log log T+O(sqrt(log log log T)))/(log\textasciicircum{}2T), while (1)/(no.I\_T)sum\_kin I\_T(tau\_k-gamma\_k)\textasciicircum{}2=(2+o(1))log log T/log\textasciicircum{}2T; the ratio is asymptotic to log log log T/(2log log T).}
\expandafter\gdef\csname CGTip@eq:computable_units\endcsname{Eq. (7.6), p. 34: < r\_k(X)\textasciicircum{}2>\_I\_T =(1)/(4pi\textasciicircum{}2)log log log T+O(sqrt(log log log T)),}
\expandafter\gdef\csname CGTip@eq:computable_height\endcsname{Eq. (7.7), p. 34: (1)/(no.I\_T)sum\_kin I\_T(tau\_k\textasciicircum{}X-gamma\_k)\textasciicircum{}2 =(log log log T+O(sqrt(log log log T)))/(log\textasciicircum{}2T),}
\expandafter\gdef\csname CGTip@thm:fixed_cutoff_step\endcsname{Theorem 7.5 (Theorem D'': the one-step clock with a fixed cutoff exponent) (p. 34): There are absolute constants T\_0 and C such that for T>= T\_0 and (log log T)\textasciicircum{}-1/2<= c<=(1)/(2), with X=T\textasciicircum{}c, < r\_k(X)\textasciicircum{}2>\_I\_T<= C(1+log log log T), < t\_k(X)\textasciicircum{}2>\_I\_T<= C(1+log log log T),; (1)/(no.I\_T)sum\_kin I\_T(tau\_k\textasciicircum{}X-gamma\_k)\textasciicircum{}2<=(C(1+log log log T))/(log\textasciicircum{}2T). In particular, for every fixed 0<c<=(1)/(2) and T>= T\_0(c)...}
\expandafter\gdef\csname CGTip@eq:fixed_cutoff_step\endcsname{Eq. (7.8), p. 34: < r\_k(X)\textasciicircum{}2>\_I\_T<= C(1+log log log T), < t\_k(X)\textasciicircum{}2>\_I\_T<= C(1+log log log T),; (1)/(no.I\_T)sum\_kin I\_T(tau\_k\textasciicircum{}X-gamma\_k)\textasciicircum{}2<=(C(1+log log log T))/(log\textasciicircum{}2T).}
\expandafter\gdef\csname CGTip@rem:fixed_cutoff_heuristic\endcsname{Remark 7.6 (What is not proved for a fixed cutoff exponent) (p. 35): Theorem 7.5 replaces the factor c\textasciicircum{}2log log T of Lemma 7.2 by 1+log log log T, the better of the two when c\textasciicircum{}2log log T>=log log log T, in particular for fixed c; but it does not give < r\_k(T\textasciicircum{}c)\textasciicircum{}2>\_I\_T=O\_c(1), which by Remark 7.9 is equivalent to bounded mean-square transport from tau\_k. The following heuristic sugges...}
\expandafter\gdef\csname CGTip@cor:local_clock\endcsname{Corollary 7.7 (A cutoff read off the clock) (p. 36): For tau\_k>e\textasciicircum{}e let X\_k\textasciicircum{}rm loc=tau\_k\textasciicircum{}1/sqrt(log log tau\_k) and tau\_k\textasciicircum{}rm loc=tau\_k-S\_X\_k\textasciicircum{}rm loc(tau\_k)/L(tau\_k). Then [ (1)/(no. I\_T\textasciicircum{}tau)sum\_kin I\_T\textasciicircum{}tau(tau\_k\textasciicircum{}rm loc-gamma\_k)\textasciicircum{}2=(log log log T+O(sqrt(log log log T)))/(log\textasciicircum{}2T). ]}
\expandafter\gdef\csname CGTip@prop:monotone\endcsname{Proposition 7.8 (The root of the corrected equation in the monotone regime) (p. 36): Let X=X(T)>=2 with M\_X<=(1)/(2)L(T), and let I'\_T be the set of kin I\_T with F\_X(T)<= k-(1)/(2)<= F\_X(2T), so that the root t\_k\textasciicircum{}Xin[T,2T] of Remark 5.7 exists. Then X<<log\textasciicircum{}2T, no.(I\_T I'\_T)<<log T, and [ (1)/(no.I\_T)sum\_kin I'\_T(t\_k\textasciicircum{}X-gamma\_k)\textasciicircum{}2>=(4)/(9L(2T)\textasciicircum{}2)((1)/(2pi\textasciicircum{}2)log log T-O(sqrt(log log T)))>>(log log T)/(log\textasciicircum{}...}
\expandafter\gdef\csname CGTip@rem:newton\endcsname{Remark 7.9 (The Newton term and the restriction on c) (p. 37): Let xi\textasciicircum{}X\_k be a point between gamma\_k and tau\_k with t\_k(X)=-S\_X'(xi\_k\textasciicircum{}X)delta\_k (mean value theorem). Since delta\_k=(r\_k(X)+S\_X(tau\_k))/L(tau\_k) by (2.12), (7.1) becomes the exact relation r\_k(X)(1+(S\_X'(xi\_k\textasciicircum{}X))/(L(tau\_k)))=u\_k(X)+chi\_k-y\_k S\_X'(xi\_k\textasciicircum{}X). The term y\_kS\_X'(xi\_k\textasciicircum{}X) is the second-order term of Newton'...}
\expandafter\gdef\csname CGTip@eq:newton_exact\endcsname{Eq. (7.9), p. 37: r\_k(X)(1+(S\_X'(xi\_k\textasciicircum{}X))/(L(tau\_k)))=u\_k(X)+chi\_k-y\_k S\_X'(xi\_k\textasciicircum{}X).}
\expandafter\gdef\csname CGTip@sec:iteration\endcsname{Section 8: Iterating the Newton-type step (p. 37)}
\expandafter\gdef\csname CGTip@subsec:recursion\endcsname{Section 8.1: The exact recursion (p. 37)}
\expandafter\gdef\csname CGTip@eq:newton_iteration\endcsname{Eq. (8.1), p. 37: x\textasciicircum{}(j)\_k=x\textasciicircum{}(j-1)\_k-(N\_0(x\textasciicircum{}(j-1)\_k)+S\_X\_j(x\textasciicircum{}(j-1)\_k)-(k-(1)/(2)))/(L(x\textasciicircum{}(j-1)\_k)), e\_j,k=x\textasciicircum{}(j)\_k-gamma\_k,}
\expandafter\gdef\csname CGTip@prop:recursion\endcsname{Proposition 8.1 (Exact recursion) (p. 37): Let j>=1 and suppose x\textasciicircum{}(j-1)\_k>2pi. Then L(x\textasciicircum{}(j-1)\_k) e\_j,k=u\_k(X\_j)+(S\_X\_j(gamma\_k)-S\_X\_j(x\textasciicircum{}(j-1)\_k))+chi\_j,k, 0<=chi\_j,k<=(e\_j-1,k\textasciicircum{}2)/(4pi min(gamma\_k,x\textasciicircum{}(j-1)\_k)) . Let varphi\_jin C\textasciicircum{}inf(R), 0<=varphi\_j<=1, be 1 on an interval containing all segments between gamma\_k and x\textasciicircum{}(j-1)\_k, kin I, and g\_j= M(varphi\_jS'\_X\_j)....}
\expandafter\gdef\csname CGTip@eq:recursion_identity\endcsname{Eq. (8.2), p. 37: L(x\textasciicircum{}(j-1)\_k) e\_j,k=u\_k(X\_j)+(S\_X\_j(gamma\_k)-S\_X\_j(x\textasciicircum{}(j-1)\_k))+chi\_j,k, 0<=chi\_j,k<=(e\_j-1,k\textasciicircum{}2)/(4pi min(gamma\_k,x\textasciicircum{}(j-1)\_k)) .}
\expandafter\gdef\csname CGTip@eq:recursion_norms\endcsname{Eq. (8.3), p. 37: ||L(x\textasciicircum{}(j-1))e\_j||\_2r,I<=||u(X\_j)||\_2r,I+2 ||e\_j-1||\_4r,I || g\_j(x\textasciicircum{}(j-1))||\_4r,I+||chi\_j||\_2r,I,; ||chi\_j||\_2r,I<=(||e\_j-1||\_4r,I\textasciicircum{}2)/(4pi min\_kin Imin(gamma\_k,x\textasciicircum{}(j-1)\_k)) .}
\expandafter\gdef\csname CGTip@lem:higher_moments\endcsname{Lemma 8.2 (p. 38): There is an absolute constant C such that for every integer r>=1, every T>= T\_0 with 2r<=kappa\_rm Flog T and every 2<= X<= T\textasciicircum{}1/(160r), with V as in Lemma 5.8 computed with Y\_2r=T\textasciicircum{}1/(160r) in place of Y\_0, [ int\_T\textasciicircum{}2TE\_X\textasciicircum{}4r dt<=(C(r\textasciicircum{}2+sqrt(r(1+V))))\textasciicircum{}4rT, < u\_k(X)\textasciicircum{}2r>\_I\_T<=(C(r\textasciicircum{}2+sqrt(r(1+V))))\textasciicircum{}2r, ] and for every r>=1...}
\expandafter\gdef\csname CGTip@thm:tails\endcsname{Theorem 8.3 (Exponential tails of the ordinal defect) (p. 39): There are absolute constants C\_8>=1, lambda\_0 and T\_0 such that, for T>= T\_0, every lambda >=lambda\_0 and every X=T\textasciicircum{}c>=2 with exp (-(lambda/C\_8)\textasciicircum{}3/2)<= c<=(1)/(C\_8sqrt(lambda)), no.kin I\_T:|u\_k(X)|>lambda<=2\textasciicircum{} 2-sqrt(lambda/C\_8) no.I\_T .}
\expandafter\gdef\csname CGTip@eq:tails_range\endcsname{Eq. (8.4), p. 39: exp (-(lambda/C\_8)\textasciicircum{}3/2)<= c<=(1)/(C\_8sqrt(lambda)),}
\expandafter\gdef\csname CGTip@eq:tails\endcsname{Eq. (8.5), p. 39: no.kin I\_T:|u\_k(X)|>lambda<=2\textasciicircum{} 2-sqrt(lambda/C\_8) no.I\_T .}
\expandafter\gdef\csname CGTip@lem:multiplicity\endcsname{Lemma 8.4 (Multiplicity) (p. 40): int\_R tilde n(u)\textasciicircum{}2 du=sum\_j,kin I\_T(b\_T-|x\textasciicircum{}(1)\_j-x\textasciicircum{}(1)\_k|)\_+<< b\_T no.I\_T .}
\expandafter\gdef\csname CGTip@eq:close_pair\endcsname{Eq. (8.6), p. 40: |y\_k+m-y\_k|>=(m-1)b .}
\expandafter\gdef\csname CGTip@lem:sampling_rho\endcsname{Lemma 8.5 (Sampling at points of bounded multiplicity) (p. 41): Let x\_k, kin I\_T, be real points of [T-2T\textasciicircum{}1/2,2T+2T\textasciicircum{}1/2] with multiplicity function tilde n(u)=no.kin I\_T:u- b\_T<= x\_k<= u satisfying int\_R tilde n\textasciicircum{}2<= k b\_Tno.I\_T for some k>=1. Let r>=1 be an integer, 2<= X<= T\textasciicircum{}1/(4r), varphi\textasciicircum{}sharp in C\textasciicircum{}inf(R) with 0<=varphi\textasciicircum{}sharp <=1, varphi\textasciicircum{}sharp=1 on [T-2T\textasciicircum{}1/2,2T+2T\textasciicircum{}1/2], varphi\textasciicircum{}sh...}
\expandafter\gdef\csname CGTip@thm:second_step\endcsname{Theorem 8.6 (Second Newton-type step) (p. 41): Let H\_0=log log T, H\_1=log H\_0, H\_2=log H\_1, X\_1=T\textasciicircum{}H\_0\textasciicircum{}-1/2, X\_2=T\textasciicircum{}H\_1\textasciicircum{}-1/2, and let x\textasciicircum{}(2)\_k be given by (8.1). Then <L(x\textasciicircum{}(1)\_k)\textasciicircum{}2e\_2,k\textasciicircum{}2>\_I\_T=(1)/(4pi\textasciicircum{}2)H\_2+O(1+sqrt(H\_2)),; (1)/(no.I\_T)sum\_kin I\_T(x\textasciicircum{}(2)\_k-gamma\_k)\textasciicircum{}2=(log log log log T+O(sqrt(log log log log T)))/(log\textasciicircum{}2T).}
\expandafter\gdef\csname CGTip@eq:second_step\endcsname{Eq. (8.7), p. 41: <L(x\textasciicircum{}(1)\_k)\textasciicircum{}2e\_2,k\textasciicircum{}2>\_I\_T=(1)/(4pi\textasciicircum{}2)H\_2+O(1+sqrt(H\_2)),; (1)/(no.I\_T)sum\_kin I\_T(x\textasciicircum{}(2)\_k-gamma\_k)\textasciicircum{}2=(log log log log T+O(sqrt(log log log log T)))/(log\textasciicircum{}2T).}
\expandafter\gdef\csname CGTip@eq:e1_fourth\endcsname{Eq. (8.8), p. 42: < e\_1,k\textasciicircum{}4>\_I\_T<<(H\_1\textasciicircum{}2)/(log\textasciicircum{}4T).}
\expandafter\gdef\csname CGTip@subsec:further_steps\endcsname{Section 8.5: All further steps: the exact identity, the local deformation, and the multiplicity (p. 42)}
\expandafter\gdef\csname CGTip@lem:second_identity\endcsname{Lemma 8.7 (The second-level identity) (p. 42): For j>=1, with y\textasciicircum{}(j)\_k=tau\_k-x\textasciicircum{}(j)\_k and chi''\_j,k=int\_x\textasciicircum{}(j-1)\_k\textasciicircum{}tau\_k(L(tau\_k)-L(y)) dy, y\textasciicircum{}(j)\_k=(1-(L(tau\_k))/(L(x\textasciicircum{}(j-1)\_k)))y\textasciicircum{}(j-1)\_k+(S\_X\_j(x\textasciicircum{}(j-1)\_k)+chi''\_j,k)/(L(x\textasciicircum{}(j-1)\_k)),; henceL(tau\_k)y\textasciicircum{}(j)\_k=S\_X\_j(x\textasciicircum{}(j-1)\_k)+ r\_j,k, with | r\_j,k|<< T\textasciicircum{}c\_j-1 if |y\textasciicircum{}(j-1)\_k|<= T\textasciicircum{}c\_j-1/2 and c\_j-1<= c\_j (true by (2.11) and i...}
\expandafter\gdef\csname CGTip@eq:Y_identity\endcsname{Eq. (8.9), p. 42: y\textasciicircum{}(j)\_k=(1-(L(tau\_k))/(L(x\textasciicircum{}(j-1)\_k)))y\textasciicircum{}(j-1)\_k+(S\_X\_j(x\textasciicircum{}(j-1)\_k)+chi''\_j,k)/(L(x\textasciicircum{}(j-1)\_k)),; henceL(tau\_k)y\textasciicircum{}(j)\_k=S\_X\_j(x\textasciicircum{}(j-1)\_k)+ r\_j,k,}
\expandafter\gdef\csname CGTip@prop:deformation\endcsname{Proposition 8.8 (Local deformation of the second clock) (p. 43): Let X\_1=T\textasciicircum{}c\_1, X\_2=T\textasciicircum{}c\_2 with c\_1=H\_0\textasciicircum{}-1/2, c\_2=H\_1\textasciicircum{}-1/2, and d\_j,k=x\textasciicircum{}(j)\_k+1-x\textasciicircum{}(j)\_k (d\_0,k=tau\_k+1-tau\_k). Then, for T>= T\_0, <|(d\_2,k-d\_0,k)/(d\_0,k)|\textasciicircum{}4>\_I\_T<< c\_2\textasciicircum{}4=(1)/(H\_1\textasciicircum{}2), (1)/(no.I\_T)no.kin I\_T:|(d\_2,k)/(d\_0,k)-1|>alpha <<(1)/(alpha\textasciicircum{}4H\_1\textasciicircum{}2)(0<alpha <=1), in particular no.kin I\_T:d\_2,k<=0<<no.I\_T/(log log log T)\textasciicircum{}2:...}
\expandafter\gdef\csname CGTip@eq:deformation\endcsname{Eq. (8.10), p. 43: <|(d\_2,k-d\_0,k)/(d\_0,k)|\textasciicircum{}4>\_I\_T<< c\_2\textasciicircum{}4=(1)/(H\_1\textasciicircum{}2), (1)/(no.I\_T)no.kin I\_T:|(d\_2,k)/(d\_0,k)-1|>alpha <<(1)/(alpha\textasciicircum{}4H\_1\textasciicircum{}2)(0<alpha <=1),}
\expandafter\gdef\csname CGTip@lem:zero_pairs\endcsname{Lemma 8.9 (Near pairs and triples of ordinates) (p. 44): Let b= b\_T, M=no.I\_T and, for R>=1 and real u, n\_ R(u)=no.lin I\_T:|gamma\_l-u|<= R b. There are absolute constants C\_1,C\_2 such that for T>= T\_0 and every R>=1, no.(k,l)in I\_T\textasciicircum{}2:|gamma\_k-gamma\_l|<= R b<= C\_1 R M,sum\_kin I\_Tn\_ R(gamma\_k)\textasciicircum{}2<= C\_2 R\textasciicircum{}2 M .}
\expandafter\gdef\csname CGTip@eq:zero_pairs\endcsname{Eq. (8.11), p. 44: no.(k,l)in I\_T\textasciicircum{}2:|gamma\_k-gamma\_l|<= R b<= C\_1 R M,sum\_kin I\_Tn\_ R(gamma\_k)\textasciicircum{}2<= C\_2 R\textasciicircum{}2 M .}
\expandafter\gdef\csname CGTip@thm:multiplicity_all\endcsname{Theorem 8.10 (Multiplicity of the corrected clock points) (p. 44): Let j>=0 and let x\textasciicircum{}(j)\_k be given by (8.1) with cutoffs X\_i=T\textasciicircum{}c\_i, c\_i<=(1)/(2), so that x\textasciicircum{}(j)\_kin[T-2T\textasciicircum{}1/2,2T+2T\textasciicircum{}1/2] for kin I\_T and T>= T\_0 (Lemma 8.7). Let tilde n\_j(u)=no.kin I\_T:u- b\_T<= x\textasciicircum{}(j)\_k<= u, and let A>=1, H>=1 be such that ||L(tau\_k)e\_j,k||\_4,I\_T<= Asqrt() H. Then, with an absolute constant C\_3, int\_R t...}
\expandafter\gdef\csname CGTip@eq:multiplicity_all\endcsname{Eq. (8.12), p. 44: int\_R tilde n\_j(u)\textasciicircum{}2 du<= C\_3 Asqrt() H b\_T no.I\_T .}
\expandafter\gdef\csname CGTip@eq:cutoff_law\endcsname{Eq. (8.13), p. 45: H\_j=2log(1)/(c\_j)(j>=1), c\_j=C\_* H\_j-1\textasciicircum{}-9/16(j>=3),}
\expandafter\gdef\csname CGTip@thm:uniform_iteration\endcsname{Theorem 8.11 (All Newton-type steps, uniformly) (p. 45): There are absolute constants C\_*, H\textasciicircum{}*, C and T\_0 such that, with the cutoffs (8.13), for every T>= T\_0 and every 1<= m<= m\textasciicircum{}*, <L(tau\_k)\textasciicircum{}2e\_m,k\textasciicircum{}2>\_I\_T=(H\_m)/(4pi\textasciicircum{}2)+vartheta\_m C(1+sqrt(H\_m)),|vartheta\_m|<=1, and ||L(tau\_k)e\_j,k||\_4,I\_T<= Csqrt(H\_j) for 1<= j<m\textasciicircum{}*. In particular, at m=m\textasciicircum{}*=log\textasciicircum{}*T+O(1), (1)/(no.I\_T)sum\_kin...}
\expandafter\gdef\csname CGTip@eq:uniform_iteration\endcsname{Eq. (8.14), p. 45: <L(tau\_k)\textasciicircum{}2e\_m,k\textasciicircum{}2>\_I\_T=(H\_m)/(4pi\textasciicircum{}2)+vartheta\_m C(1+sqrt(H\_m)),|vartheta\_m|<=1,}
\expandafter\gdef\csname CGTip@eq:final_clock\endcsname{Eq. (8.15), p. 45: (1)/(no.I\_T)sum\_kin I\_T(x\textasciicircum{}(m\textasciicircum{}*)\_k-gamma\_k)\textasciicircum{}2<=(C')/(log\textasciicircum{}2T):}
\expandafter\gdef\csname CGTip@eq:uniform_recursion\endcsname{Eq. (8.16), p. 46: A\_j(r)<= C(sqrt(r)+(r\textasciicircum{}2)/(sqrt(H\_j)))+(C\_5sqrt(r))/(sqrt(H\_j))max (A\_j-1(2r),A\_j-1(2r)\textasciicircum{}9/8),}
\expandafter\gdef\csname CGTip@eq:induction_claim\endcsname{Eq. (8.17), p. 46: A\_j(r)<= B\_* r\textasciicircum{}2(0<= j<= m-1, 1<= r<=2\textasciicircum{}m-j).}
\expandafter\gdef\csname CGTip@cor:third_step\endcsname{Corollary 8.12 (Third step) (p. 47): With c\_3=C\_* H\_2\textasciicircum{}-9/16 and H\_3=(9)/(8)log H\_2+2log(1/C\_*), [ <L(tau\_k)\textasciicircum{}2e\_3,k\textasciicircum{}2>\_I\_T=(H\_3)/(4pi\textasciicircum{}2)+O(1+sqrt(H\_3)),(1)/(no.I\_T)sum\_kin I\_T(x\textasciicircum{}(3)\_k-gamma\_k)\textasciicircum{}2=(H\_3+O(sqrt(H\_3)))/(log\textasciicircum{}2T), ] and the multiplicity of the second clock satisfies int tilde n\_2\textasciicircum{}2llsqrt(H\_2) b\_Tno.I\_T.}
\expandafter\gdef\csname CGTip@cor:clock_tails\endcsname{Corollary 8.13 (Tails of the corrected clocks) (p. 47): With the cutoffs (8.13), for T>= T\_0, every 1<= j<m\textasciicircum{}*, every integer 1<= q<=2\textasciicircum{}m\textasciicircum{}*-j and every lambda>0, no.kin I\_T:L(tau\_k)|e\_j,k|>lambdasqrt(H\_j)<=((B\_* q\textasciicircum{}2)/(lambda))\textasciicircum{}2qno.I\_T . For fixed j and q the right side is <<\_qlambda\textasciicircum{}-2qno.I\_T for T>= T\_0(j,q), since 2\textasciicircum{}m\textasciicircum{}*-j->inf: for every fixed q, and T large in terms of q, th...}
\expandafter\gdef\csname CGTip@eq:clock_tails\endcsname{Eq. (8.18), p. 47: no.kin I\_T:L(tau\_k)|e\_j,k|>lambdasqrt(H\_j)<=((B\_* q\textasciicircum{}2)/(lambda))\textasciicircum{}2qno.I\_T .}
\expandafter\gdef\csname CGTip@lem:sorted\endcsname{Lemma 8.14 (Sorting the final points) (p. 48): Let gamma\_1<=...<=gamma\_ M, let z\_1,...,z\_ M be real, and let z\textasciicircum{}increasing to be the same values in nondecreasing order. Then sum\_k|z\textasciicircum{}increasing to\_k-gamma\_k|\textasciicircum{}q<=sum\_k|z\_k-gamma\_k|\textasciicircum{}q for every q>=1, and z\textasciicircum{}increasing to\_k=z\_k as multisets. In particular, sorting the points x\textasciicircum{}(j)\_k, kin I, of one level does not increa...}
\expandafter\gdef\csname CGTip@cor:sorted_clock\endcsname{Corollary 8.15 (A nondecreasing clock built without ordinates) (p. 48): Let I\_T\textasciicircum{}tau=k:T<=tau\_k<2T, a set of consecutive indices defined by the clock alone. For each T, with the cutoffs and the number of steps chosen at height T, the points x\textasciicircum{}(m\textasciicircum{}*)\_k of Theorem 8.11 are defined by (8.1) on the block J\_T=k:T/2<=tau\_k<=3T, which contains I\_T u I\_T\textasciicircum{}tau for T>= T\_0; let x\textasciicircum{}increasing to\_k, ki...}
\expandafter\gdef\csname CGTip@sec:regularized\endcsname{Section 9: Theorem D': a regularized clock with a fixed cutoff (p. 49)}
\expandafter\gdef\csname CGTip@def:regularized\endcsname{Definition 9.1 (Regularized clock) (p. 49): hat F=hat F\_T,X is the element of K\_T that minimizes int\_ J\_T|P-F\_X|\textasciicircum{}4 dt over Pin K\_T. For every integer k with hat F((T)/(2))<k-(1)/(2)<hat F(3T), hat x\_k=hat x\_k(T,X)in J\_T is the solution of hat F(hat x\_k)=k-(1)/(2), and hat u\_k=k-(1)/(2)-hat F(gamma\_k) when gamma\_kin J\_T.}
\expandafter\gdef\csname CGTip@lem:regularized_def\endcsname{Lemma 9.2 (The regularized clock is well defined) (p. 49): For T>= T\_0 and X>=2 the minimizer hat F exists and is unique; hat F is strictly increasing with (1)/(2)Lambda\_T<=hat F'<=2Lambda\_T on each mesh interval; each hat x\_k is unique, obtained by locating the mesh interval where hat F crosses k-(1)/(2) and interpolating linearly; and (pi)/(log T)<=hat x\_k+1-hat x\_k<=(4pi...}
\expandafter\gdef\csname CGTip@eq:regularized_gaps\endcsname{Eq. (9.1), p. 49: (pi)/(log T)<=hat x\_k+1-hat x\_k<=(4pi)/(log T)}
\expandafter\gdef\csname CGTip@lem:identity_lipschitz\endcsname{Lemma 9.3 (Theorem A for Lipschitz models) (p. 50): Theorems 4.1 and 4.2 hold for every real Lipschitz function F on an open interval containing [gamma\_k\_1,gamma\_k\_2], with F' its derivative almost everywhere; in particular for the functions of K\_T.}
\expandafter\gdef\csname CGTip@lem:slope_envelope\endcsname{Lemma 9.4 (A slope-constrained model of the counting function) (p. 50): There are absolute constants T\_0 and C such that for T>= T\_0 there is P\_Nin K\_T with int\_ J\_T|P\_N-N|\textasciicircum{}4 dt<= CT.}
\expandafter\gdef\csname CGTip@lem:block_fourth\endcsname{Lemma 9.5 (The prime model on the block) (p. 51): There are absolute constants T\_0 and C such that for T>= T\_0 and 2<= Xlesqrt(T/2), with c=log X/log T, [ int\_ J\_T|F\_X-N|\textasciicircum{}4 dt<= CT(1+log(1/c))\textasciicircum{}2 . ]}
\expandafter\gdef\csname CGTip@thm:regularized_clock\endcsname{Theorem 9.6 (Theorem D') (p. 51): There are absolute constants T\_0 and C such that for T>= T\_0 and 2<= Xlesqrt(T/2), with c=log X/log T and hat F, hat x\_k, hat u\_k as in Definition 9.1, and with I either of the index sets I\_T and I\_T\textasciicircum{}tau=k:T<=tau\_k<2T: item hat x\_k is defined for every kin I\_T u I\_T\textasciicircum{}tau, and these points are increasing in k with gap...}
\expandafter\gdef\csname CGTip@eq:regularized_main\endcsname{Eq. (9.2), p. 51: <|hat x\_k-gamma\_k|\textasciicircum{}4>\_ I<=(C(1+log(1/c))\textasciicircum{}2)/(log\textasciicircum{}4T),<|hat x\_k-gamma\_k|\textasciicircum{}2>\_ I<=(C(1+log(1/c)))/(log\textasciicircum{}2T),}
\expandafter\gdef\csname CGTip@eq:regularized_tails\endcsname{Eq. (9.3), p. 51: (no.kin I:|hat x\_k-gamma\_k|>lambda *2pi/log T)/(no. I)<=(C(1+log(1/c))\textasciicircum{}2)/((2pi lambda)\textasciicircum{}4);}
\expandafter\gdef\csname CGTip@eq:regularized_close\endcsname{Eq. (9.4), p. 52: ||hat F-N||\_4<=||hat F-F\_X||\_4+||F\_X-N||\_4<=||P\_N-F\_X||\_4+||F\_X-N||\_4<=||P\_N-N||\_4+2||F\_X-N||\_4,}
\expandafter\gdef\csname CGTip@subsec:tolerance\endcsname{Section 9.4: A certified tolerance for the convex program (p. 53)}
\expandafter\gdef\csname CGTip@prop:tolerance\endcsname{Proposition 9.7 (Approximate minimizers and a dual certificate) (p. 53): Let T, X, c be as in Theorem 9.6 and 0<=eta\_rm reg<=1. item If P\_ vin K\_T satisfies E(v)<= E\_*+eta\_rm regT, and hat x\_k\textasciicircum{}(v) are its level points, P\_ v(hat x\_k\textasciicircum{}(v))=k-(1)/(2), then (a)--(d) of Theorem 9.6 hold for these points, with (1+log(1/c))\textasciicircum{}2 replaced by (1+log(1/c))\textasciicircum{}2+eta\_rm reg in the fourth-moment bounds and...}
\expandafter\gdef\csname CGTip@rem:regularized\endcsname{Remark 9.8 (p. 54): (a) Compared with Theorem 8.11 and Corollary 8.15, Theorem 9.6 reaches the scale of one mean spacing with a fixed cutoff exponent and without Newton iteration, in fourth mean, with increasing and separated points, on the clock window I\_T\textasciicircum{}tau as well as on I\_T; the dependence on c is explicit and parallel to Theorem...}
\expandafter\gdef\csname CGTip@sec:envelope\endcsname{Section 10: Theorem E: the pointwise envelope (p. 54)}
\expandafter\gdef\csname CGTip@prop:envelope_equivalence\endcsname{Proposition 10.1 (The envelope of the clock is the envelope of S) (p. 54): (a) limsup\_k->inf|tau\_k-gamma\_k|=2pi a. (b) If X\_k>=2 and sigma\_X\_k=o(log tau\_k), in particular if X\_k<=(log tau\_k)\textasciicircum{}2-varepsilon, then limsup\_k->inf|tau\_k\textasciicircum{}X\_k-gamma\_k|=2pi a.}
\expandafter\gdef\csname CGTip@thm:envelope\endcsname{Theorem 10.2 (Theorem E) (p. 55): Let e be continuous on [(1)/(2),1] with e(1)=0, such that for every varepsilon>0, zeta(sigma+it)<<\_varepsilon|t|\textasciicircum{} e(sigma)+varepsilon uniformly for (1)/(2)<=sigma <=1 and |t|>=2, and put e(sigma)=(1)/(2)-sigma+ e(1-sigma) for 0<=sigma<(1)/(2). Then a<=(1)/(2pi log 2)int\_0\textasciicircum{}1(e(sigma))/(sqrt(sigma(2-sigma))) dsigma ....}
\expandafter\gdef\csname CGTip@eq:envelope_formula\endcsname{Eq. (10.1), p. 55: a<=(1)/(2pi log 2)int\_0\textasciicircum{}1(e(sigma))/(sqrt(sigma(2-sigma))) dsigma .}
\expandafter\gdef\csname CGTip@eq:envelope_value\endcsname{Eq. (10.2), p. 55: C\_inf=0.0968151142790...,limsup\_k->inf|tau\_k-gamma\_k|=2pi a<=2pi C\_inf=0.608307303... ;}
\expandafter\gdef\csname CGTip@eq:envelope_value_hull\endcsname{Eq. (10.3), p. 55: C\_ P=0.0966530245541...,limsup\_k->inf|tau\_k-gamma\_k|<=2pi C\_ P=0.6072888637...<0.6072889 .}
\expandafter\gdef\csname CGTip@rem:envelope_inputs\endcsname{Remark 10.3 (p. 57): (a) Trudgian [ref, S 5 of the preprint] obtained a<=0.1013 from Huxleytextquoteright s exponent (32)/(205) at sigma=(1)/(2) by the same method; Theorem 10.2 replaces the exponents and lets the parameters of the disc tend to their limits. With the explicit exponents used by Bellotti and Wong ((27)/(164) at sigma=(1)/...}
\expandafter\gdef\csname CGTip@sec:simple_zeros\endcsname{Section 11: Theorem F: simple zeros on the critical line from the defect (p. 57)}
\expandafter\gdef\csname CGTip@eq:Q_def\endcsname{Eq. (11.1), p. 57: Q\_2j(T,X)=(1)/(no.I\_T)int\_T\textasciicircum{}2T(N(t)-F\_X(t))\textasciicircum{}2jF\_X'(t) dt .}
\expandafter\gdef\csname CGTip@lem:mult_ordinate\endcsname{Lemma 11.1 (Multiplicity of an ordinate) (p. 58): item If an ordinate gamma occurs with multiplicity m>=2 in the list (gamma\_k), at the indices k,...,k+m-1, then u\_k+i=u\_k+i for 0<= i<m, and for every j>=1 [ sum\_i=0\textasciicircum{}m-1(2u\_k+i)\textasciicircum{}2j>= m . ] item If gamma occurs with multiplicity 1, the zero with ordinate gamma is simple and lies on the critical line.}
\expandafter\gdef\csname CGTip@thm:simple_certificate\endcsname{Theorem 11.2 (Theorem F) (p. 58): For every real function F, every j>=1 and every T>0, N\_0\textasciicircum{}s(T,2T)>=no.I\_T-sum\_kin I\_T(2u\_k)\textasciicircum{}2j.}
\expandafter\gdef\csname CGTip@eq:simple_certificate\endcsname{Eq. (11.2), p. 58: N\_0\textasciicircum{}s(T,2T)>=no.I\_T-sum\_kin I\_T(2u\_k)\textasciicircum{}2j.}
\expandafter\gdef\csname CGTip@cor:simple_Q\endcsname{Corollary 11.3 (p. 58): Let j>=1 and 0<theta<(2)/(2j+1) be fixed. For F=F\_X with 2<= X<= T\textasciicircum{}theta, as T->inf, (N\_0\textasciicircum{}s(T,2T))/(no.I\_T)>=1-4\textasciicircum{}j< u\_k(X)\textasciicircum{}2j>\_I\_T,< u\_k(X)\textasciicircum{}2>\_I\_T=Q\_2(T,X)-(1)/(12)+o(1), so that N\_0\textasciicircum{}s(T,2T)/no.I\_T >=(4)/(3)-4Q\_2(T,X)-o(1) for X<= T\textasciicircum{}1/2, and N\_0\textasciicircum{}s(T,2T)/no.I\_T >=1-16Q\_4(T,X)+8Q\_2(T,X)-(7)/(15)-o(1) for X<= T\textasciicircum{}1/3.}
\expandafter\gdef\csname CGTip@eq:simple_Q\endcsname{Eq. (11.3), p. 58: (N\_0\textasciicircum{}s(T,2T))/(no.I\_T)>=1-4\textasciicircum{}j< u\_k(X)\textasciicircum{}2j>\_I\_T,< u\_k(X)\textasciicircum{}2>\_I\_T=Q\_2(T,X)-(1)/(12)+o(1),}
\expandafter\gdef\csname CGTip@rem:simple_threshold\endcsname{Remark 11.4 (What the certificate needs) (p. 59): The cross term (1)/(no.I\_T)int\_T\textasciicircum{}2T(N-F\_X)\textasciicircum{}2S\_X' dt, which comes from dF\_X=(L+S\_X') dt, is part of Q\_2; no hypothesis on the zeros enters (11.3). The first form gives a positive proportion as soon as Q\_2(T,T\textasciicircum{}c)<=(1)/(3)-eta for one fixed c<=(1)/(2), eta>0 and all large T, and the proportion (2)/(3) when Q\_2<=(1)/(6)....}
\expandafter\gdef\csname CGTip@sec:euler_product\endcsname{Section 12: Theorem G: the finite Euler product in the zeta function of the corrected roots (p. 59)}
\expandafter\gdef\csname CGTip@eq:zeta_kappa\endcsname{Eq. (12.1), p. 59: zeta\_kappa(s)=(1)/(G(s))prod\_k>=1(1+(s(s-1))/((1)/(4)+kappa\_k\textasciicircum{}2)), G(s)=s(s-1)pi\textasciicircum{}-s/2Gamma ((s)/(2)).}
\expandafter\gdef\csname CGTip@eq:f_z\endcsname{Eq. (12.2), p. 59: 1+(s(s-1))/((1)/(4)+t\textasciicircum{}2)=(t\textasciicircum{}2+z\textasciicircum{}2)/(t\textasciicircum{}2+(1)/(4)), f\_z(t)=log(t\textasciicircum{}2+z\textasciicircum{}2)/(t\textasciicircum{}2+(1)/(4))(Re z>0, tin R),}
\expandafter\gdef\csname CGTip@lem:fourier_log\endcsname{Lemma 12.1 (p. 59): For Re z>0 and omega>0, [ int\_0\textasciicircum{}inf f\_z(t)cos(omega t) dt=(pi)/(omega)(e\textasciicircum{}-omega/2-e\textasciicircum{}-omega z). ]}
\expandafter\gdef\csname CGTip@lem:midpoint_rule\endcsname{Lemma 12.2 (The midpoint rule along a monotone model) (p. 60): Let Fin C\textasciicircum{}1[t\_*,inf) with F'>0 on [t\_*,inf), F(t)->inf and F(t\_*)=v\_*in Z, and let rho\_k be the root of F(t)=k-(1)/(2), k>v\_*. Let phi in C\textasciicircum{}1[t\_*,inf) with phi ->0 at inf and int\_t\_*\textasciicircum{}inf(|phi| F'+|phi'|) dt<inf. Then, with the 1-periodic sawtooth psi(v)=floor(v+(1)/(2))-v, [ sum\_k>v\_*phi(rho\_k)=int\_t\_*\textasciicircum{}inf phi(t) F'...}
\expandafter\gdef\csname CGTip@thm:euler_product\endcsname{Theorem 12.3 (Theorem G) (p. 60): Let X>=2 be fixed and M\_X the bound for |S\_X'| in (2.11). Let t\_1>2pi with L(t\_1)>=2M\_X, so that F\_X'>=(1)/(2)L>0 on [t\_1,inf); let v\_0=ceil(F\_X(t\_1)), t\_0=F\_X\textasciicircum{}-1(v\_0), k\_0=v\_0+1, and for k>= k\_0 let kappa\_k be the unique root of F\_X(t)=k-(1)/(2) in [t\_0,inf); for 1<= k<k\_0 let kappa\_k be any nonzero reals, for inst...}
\expandafter\gdef\csname CGTip@eq:euler_product\endcsname{Eq. (12.3), p. 60: sum\_k>=1(f\_z(kappa\_k)-f\_z(tau\_k))=log P\_X(s)-log P\_X(1)+H\_X(s),}
\expandafter\gdef\csname CGTip@rem:euler_product\endcsname{Remark 12.4 (p. 61): The zeros of zeta\_kappa are on the critical line by construction, and zeta\_tau, P\_X and e\textasciicircum{}H\_X are free of zeta: Theorem G separates, in the zeta function of the corrected roots, the part supplied by the primes, which is exactly the finite Euler product P\_X, from the part supplied by the clock. Consequently, in Re s>...}
\expandafter\gdef\csname CGTip@sec:conclusions\endcsname{Section 13: Conclusions (p. 61)}
\makeatother
